\documentclass[a4paper,10pt]{article}
\usepackage[czech]{babel}
\usepackage[T1]{fontenc}
\usepackage[utf8]{inputenc}
\usepackage{amsmath,amssymb}
\usepackage[margin=18mm]{geometry}
\setlength{\parindent}{0pt}

\begin{document}
\section*{Komplexní čísla A/B -- řešení}
Podklad pro učitele k testu \texttt{02\_test\_KomplexniCislaZaverecny\_AB.tex}.

Zdroj všech deseti zadání: přímé zadání učitele v této konverzaci,
ve zprávě začínající „Vytvoř test A/B na komplexní čísla podle následujících příkladů.“
Úloha A1 až A5 odpovídá položce 1 až 5 pod nadpisem „Varianta A“;
obdobně B1 až B5 položkám pod nadpisem „Varianta B“.
Zadání ani přiřazení k variantám nebyla změněna.

Zařazení: 3.~ročník. Tematické plány:\\
\texttt{3\_rocnik/tematicke\_plany/SPS.pdf} a
\texttt{3\_rocnik/tematicke\_plany/GYM.pdf}, s.~1.
Oba plány řadí komplexní čísla na začátek třetího ročníku;
početní operace a Moivreovu větu do září, kvadratické rovnice do října.
Zářijová dotace činí 9 hodin pro SPŠ a 12 hodin pro gymnázium.

\subsection*{Varianta A}
\begin{enumerate}
\item \((2+3i)^2\).
Použijeme vzorec pro druhou mocninu dvojčlenu a vztah $i^2=-1$:
\[
(2+3i)^2=4+12i+9i^2=4+12i-9=-5+12i.
\]
Nezávislá kontrola přímým násobením:
$(2+3i)(2+3i)=4+6i+6i-9=-5+12i$.

\item \(\left|\frac{2+3i}{1-i}\right|\).
Protože $1-i\neq0$, platí
\[
\left|\frac{2+3i}{1-i}\right|
=\frac{|2+3i|}{|1-i|}
=\frac{\sqrt{2^2+3^2}}{\sqrt{1^2+(-1)^2}}
=\frac{\sqrt{13}}{\sqrt2}=\frac{\sqrt{26}}2.
\]
Kontrola jiným postupem:
\[
\frac{2+3i}{1-i}=\frac{(2+3i)(1+i)}2=-\frac12+\frac52i,
\qquad
\sqrt{\left(-\frac12\right)^2+\left(\frac52\right)^2}=\frac{\sqrt{26}}2.
\]

\item \(x^2-5x+5=0\), řešte v \(\mathbb C\).
Diskriminant je $D=(-5)^2-4\cdot1\cdot5=5>0$, tedy
\[
x_{1,2}=\frac{5\pm\sqrt5}{2},
\qquad
K=\left\{\frac{5-\sqrt5}{2};\frac{5+\sqrt5}{2}\right\}.
\]
Oba kořeny jsou reálné, a tedy patří i do $\mathbb C$.
Zkouška pro obě znaménka:
\[
\left(\frac{5\pm\sqrt5}{2}\right)^2
-5\frac{5\pm\sqrt5}{2}+5
=\frac{30\pm10\sqrt5-50\mp10\sqrt5+20}{4}=0.
\]
Jde o dva různé kořeny kvadratické rovnice; jiné kořeny neexistují.

\item \(z_1=1+i\),
\(z_2=4\left(\cos\frac{5\pi}{6}+i\sin\frac{5\pi}{6}\right)\).
Určete \(\frac{z_2}{z_1}\) v algebraickém tvaru.

Nejprve $z_2=4(-\frac{\sqrt3}{2}+\frac12i)=-2\sqrt3+2i$.
Protože $z_1\neq0$, můžeme dělit:
\[
\frac{z_2}{z_1}
=\frac{(-2\sqrt3+2i)(1-i)}{(1+i)(1-i)}
=\frac{2-2\sqrt3+(2+2\sqrt3)i}{2}
=(1-\sqrt3)+(1+\sqrt3)i.
\]
Kontrola zpětným násobením:
$[(1-\sqrt3)+(1+\sqrt3)i](1+i)=-2\sqrt3+2i=z_2$.

\item \(z=\sqrt2\left(\cos\frac{3\pi}{4}+i\sin\frac{3\pi}{4}\right)\).
Pomocí Moivreovy věty vypočtěte \(z^4\) a výsledek zapište v algebraickém tvaru.
\[
z^4=(\sqrt2)^4\left(\cos\left(4\cdot\frac{3\pi}{4}\right)
+i\sin\left(4\cdot\frac{3\pi}{4}\right)\right)
=4(\cos3\pi+i\sin3\pi)=-4.
\]
Kontrola v algebraickém tvaru: $z=-1+i$, $z^2=-2i$, tedy $z^4=(-2i)^2=-4$.
\end{enumerate}

\newpage
\section*{Komplexní čísla -- varianta B}
Řešení k testu \texttt{02\_test\_KomplexniCislaZaverecny\_AB.tex}.
\begin{enumerate}
\item \((4-i)(2+5i)\).
\[
(4-i)(2+5i)=8+20i-2i-5i^2=13+18i.
\]
Kontrola odděleným výpočtem složek součinu:
$\operatorname{Re}=4\cdot2-(-1)\cdot5=13$,
$\operatorname{Im}=4\cdot5+(-1)\cdot2=18$.

\item \(\left|\frac{4-i}{2+2i}\right|\).
Protože $2+2i\neq0$, platí
\[
\left|\frac{4-i}{2+2i}\right|
=\frac{|4-i|}{|2+2i|}
=\frac{\sqrt{4^2+(-1)^2}}{\sqrt{2^2+2^2}}
=\frac{\sqrt{17}}{\sqrt8}=\frac{\sqrt{34}}4.
\]
Kontrola jiným postupem:
\[
\frac{4-i}{2+2i}=\frac{(4-i)(2-2i)}8=\frac34-\frac54i,
\qquad
\sqrt{\left(\frac34\right)^2+\left(-\frac54\right)^2}=\frac{\sqrt{34}}4.
\]

\item \(x^2+4x+5=0\), řešte v \(\mathbb C\).
Doplníme na druhou mocninu:
\[
x^2+4x+5=0 \iff (x+2)^2=-1 \iff x+2=\pm i.
\]
Proto $K=\{-2-i;-2+i\}$.
Zkouška pro obě znaménka:
\[
(-2\pm i)^2+4(-2\pm i)+5
=3\mp4i-8\pm4i+5=0.
\]
Jde o dva různé kořeny kvadratické rovnice; jiné kořeny neexistují.

\item \(z_1=1-i\),
\(z_2=6\left(\cos\frac{2\pi}{3}+i\sin\frac{2\pi}{3}\right)\).
Určete \(\frac{z_2}{z_1}\) v algebraickém tvaru.

Nejprve $z_2=6(-\frac12+\frac{\sqrt3}{2}i)=-3+3\sqrt3i$.
Protože $z_1\neq0$, můžeme dělit:
\[
\frac{z_2}{z_1}
=\frac{(-3+3\sqrt3i)(1+i)}{(1-i)(1+i)}
=-\frac{3(1+\sqrt3)}2+\frac{3(\sqrt3-1)}2i.
\]
Kontrola zpětným násobením:
\[
\left[-\frac{3(1+\sqrt3)}2+\frac{3(\sqrt3-1)}2i\right](1-i)
=-3+3\sqrt3i=z_2.
\]

\item \(z=2\left(\cos\frac{\pi}{6}+i\sin\frac{\pi}{6}\right)\).
Pomocí Moivreovy věty vypočtěte \(z^3\) a výsledek zapište v algebraickém tvaru.
\[
z^3=2^3\left(\cos\left(3\cdot\frac{\pi}{6}\right)
+i\sin\left(3\cdot\frac{\pi}{6}\right)\right)
=8\left(\cos\frac\pi2+i\sin\frac\pi2\right)=8i.
\]
Kontrola v algebraickém tvaru: $z=\sqrt3+i$, $z^2=2+2\sqrt3i$, tedy
\[
z^3=(2+2\sqrt3i)(\sqrt3+i)
=2\sqrt3+2i+6i-2\sqrt3=8i.
\]
\end{enumerate}

\subsection*{Kontrola a srovnatelnost variant}
Všechny podíly mají nenulové jmenovatele. U absolutních hodnot je výsledek
nezáporný. Obě kvadratické rovnice mají právě dva uvedené kořeny v $\mathbb C$.
Početní výsledky byly navíc nezávisle ověřeny přesnou aritmetikou;
goniometrická zadání také numerickým výpočtem z původních úhlů.

Varianty ověřují stejné okruhy a mají přibližně srovnatelný rozsah.
V A1 je druhá mocnina, v B1 součin; v A3 jsou dva reálné iracionální kořeny,
v B3 dva nereálné kořeny. Tyto rozdíly vycházejí přímo ze zadání učitele.
Orientační časy úloh jsou 2--3, 3--4, 3--5, 5--7 a 3--5 minut;
celkem přibližně 20--30 minut včetně kontroly. Jde o odhad, nikoliv pevně
stanovenou délku testu nebo údaj ověřený ve třídě.
\end{document}
